\honest{Distinct Configurations}{Eliezer Yudkowsky}{2008}

The last post taught two lessons: which configurations are the same is a physical fact, and a configuration is about all the particles, in the end ``all the particles… everywhere.'' \nb{That a single quantum state describes the whole universe is Everett's view. Bohr rejected it.}

Now I put a ``sensitive thingy,'' S, on one path of the interferometer. S changes from No to Yes when a photon passes. Every configuration now includes S, so the two flows that used to cancel at Detector 1 end in different configurations, one with S at No and one with S at Yes. Each of the four final configurations gets a quarter, and the photon reaches each detector half the time. \nb{The arithmetic is right. The figure still draws the path to Detector 1 as a dashed line, the previous post's sign for no photons.}

This is true if nobody cares about S, if nobody looks, if nobody knows S exists, and if S's state flies off into space where no one could ever learn it. It fails only if S is moved by less than its original spread. \nb{This is standard physics, and the condition is right: interference is lost in degrees, according to how well S's two states can be told apart.} The summary: the configurations are distinct ``so long as at least one particle in the universe anywhere is in a different position.'' \nb{The summary leaves out the condition just given.} ``Configurations are not belief states.'' \nb{The experiment shows that what anyone knows does not matter. It does not show whether amplitudes are physical; every interpretation predicts this result.}

Then history. An early scientist, I imagine, sees the interference vanish when a sensor is added and decides that the photon ``doesn't want me looking at it too closely.'' ``STOP THE PRESSES! MIND IS FUNDAMENTAL AFTER ALL!'' You can still read this ``In physics textbooks,'' even now, ``when a majority of theoretical physicists know better.'' \nb{No textbook is named. In one small survey of foundations researchers, 6 per cent gave the observer a special physical role.}

The clues were there: an unread sensor does the same, and so does a single stray particle. I ask whether you could have done better than John von Neumann, and call it ``a little embarrassing'' that the view survived after the theory predicted that any nudged particle would do. \nb{Von Neumann had shown that the predictions do not depend on where the line between system and observer is drawn. On his account, too, an unread sensor removes the interference; what he tied to consciousness was the appearance of a single outcome, which these clues do not test.}

The ``Common Wisdom'' was that configurations were possibilities, amplitudes were partial information, and ``conscious observation made quantumness go away.'' \nb{That merges several views. The consciousness part was von Neumann's and Wigner's; Bohr treated the interaction with the apparatus as objective.}
